\documentclass[11pt]{article}
\usepackage[margin=0.9in]{geometry}
\usepackage{amsmath,amssymb,amsthm,graphicx,microtype,enumitem}
\usepackage[hidelinks]{hyperref}

\newtheorem{theorem}{Theorem}
\newtheorem{lemma}[theorem]{Lemma}
\newtheorem{corollary}[theorem]{Corollary}
\theoremstyle{remark}
\newtheorem{remark}[theorem]{Remark}
\newcommand{\Ghat}{\widehat G}
\newcommand{\Epp}{\mathcal E_{\mathrm{KD}+}^{\mathrm{pure}}}
\newcommand{\VKD}{V_{\mathrm{KDr}}}
\newcommand{\HS}{\mathcal T_2(G)}
\newcommand{\X}{\mathfrak X}
\newcommand{\one}{\mathbf 1}

\title{KD-real observables are generated by KD-positive pure states}
\author{Candidate full result for arXiv:2507.23628}
\date{13 August 2026}

\begin{document}
\maketitle

\begin{center}
\fbox{\parbox{0.92\textwidth}{\textbf{Status.}
Candidate full result, likely valid.  We prove Conjecture (11) of the source
for every second-countable locally compact abelian group.  The separate
conjecture about the convex hull of KD-positive states is not claimed.}}
\end{center}

\section{Source conjecture}

The source is Mat\'eo Spriet, \emph{Characterizing the Kirkwood--Dirac
positivity on second countable LCA groups}, arXiv:2507.23628 (current PDF
dated 23 February 2026) \cite{Spriet}.  Its Theorem 1.3 proves two structural
equalities for compact connected groups.  The first is then conjectured to
hold for every second-countable LCA (SCLCA) group.  Figure~\ref{fig:conj}
reproduces the exact passage.

\begin{figure}[ht]
\centering
\includegraphics[width=0.94\textwidth]{figures/conjecture_crop.png}
\caption{The two conjectures on page 7 of the official arXiv PDF.}
\label{fig:conj}
\end{figure}

\noindent\textbf{Transcription of Conjecture (11).}
If $\VKD$ is the real vector space of self-adjoint Hilbert--Schmidt
observables with real Kirkwood--Dirac distribution, and $\Epp$ is the set of
KD-positive pure states, then
\[
 \VKD=\overline{\operatorname{span}_{\mathbb R}(\Epp)}^{\,\|\cdot\|_2}
 \qquad\text{for every SCLCA group }G.                         \tag{1.1}
\]

\section{Setup from the source}

Fix dual Haar measures on $G$ and $\Ghat$.  Let
\[
 U(u,\eta)=M_\eta T_u,\qquad
 \X_A^0(u,\eta)=\operatorname{Tr}(A U(u,\eta)).
\]
The source proves that $A\mapsto\X_A^0$ extends to a unitary map
\[
 \HS\longrightarrow L^2(G\times\Ghat).                       \tag{2.1}
\]
Its Lemma 4.1 says that a self-adjoint $A\in\HS$ belongs to $\VKD$ exactly
when $\X_A^0$ is essentially supported on
\[
 N=\{(u,\eta)\in G\times\Ghat:\eta(u)=1\}.                   \tag{2.2}
\]

For a compact open subgroup $H\leq G$, write
\[
 H^\perp=\{\eta\in\Ghat:\eta(h)=1\ (h\in H)\},\qquad
 e_H=\mu_G(H)^{-1/2}\one_H,
\]
and $P_H=|e_H\rangle\langle e_H|$.  The source's Theorem 4.1 classifies
all elements of $\Epp$: they are exactly
\[
 P_{H,a,\alpha}=U(a,\alpha)P_HU(a,\alpha)^*,                 \tag{2.3}
\]
where $H$ is compact open and
$(a,\alpha)\in G/H\times\Ghat/H^\perp$.

\section{Proof intuition}

The support condition (2.2) converts the operator conjecture into a density
problem on phase space.  For each compact open $H$, the rectangle
$H\times H^\perp$ lies in $N$.  The characteristic function of $P_H$ is its
indicator.  Weyl translation multiplies that indicator by a character, and
the permitted translation parameters give \emph{every} character of the
compact group $H\times H^\perp$.  Thus the pure states associated with one
$H$ already span the full $L^2$ space on that rectangle.

It remains to cover $N$.  A fiber of $N$ over $u$ has positive Haar measure
only when the cyclic closure of $u$ is compact; the symmetric assertion holds
for the $\eta$ coordinate.  For almost every $(u,\eta)\in N$, both cyclic
closures are therefore compact.  Starting with any compact open subgroup,
one can enlarge a smaller compact open subgroup by the cyclic closure of $u$
while staying inside $\ker\eta$.  This puts $(u,\eta)$ in some
$H\times H^\perp$.  Second countability reduces the cover to countably many
rectangles, completing the density argument.

\section{The phase-space covering}

Call $x$ in an LCA group \emph{topologically periodic} when
$\overline{\langle x\rangle}$ is compact.

\begin{lemma}[rectangle cover]\label{lem:cover}
Suppose $G$ has a compact open subgroup.  Then, up to a product-Haar-null
set,
\[
 N=\bigcup_{H\leq G\,\text{compact open}} H\times H^\perp.   \tag{4.1}
\]
Moreover, a countable subfamily of the rectangles in (4.1) has the same
union up to a null set.
\end{lemma}

\begin{proof}
For fixed $u\in G$, the vertical fiber of $N$ is
\[
 N_u=\{\eta:\eta(u)=1\}
     =\overline{\langle u\rangle}^{\,\perp}.
\]
A closed subgroup of an LCA group has positive Haar measure exactly when it
is open.  By Pontryagin duality,
$\overline{\langle u\rangle}^{\perp}$ is open exactly when
$\overline{\langle u\rangle}$ is compact.  Tonelli's theorem therefore shows
that, for almost every $(u,\eta)\in N$, $u$ is topologically periodic.
Applying the same argument to the horizontal fibers shows that $\eta$ is
also topologically periodic almost everywhere.

Consider such a pair $(u,\eta)\in N$.  Put
$C=\overline{\langle u\rangle}\leq G$ and
$D=\overline{\langle\eta\rangle}\leq\Ghat$; both are compact.  Since
$\eta(u)=1$, continuity gives $C\subseteq\ker\eta$.  Compactness of $D$
implies that $D^\perp$ is open in $G$.  Fix one compact open $K_0\leq G$ and
let
\[
 K=K_0\cap D^\perp,\qquad H=C+K.
\]
Then $K$ is compact open, so $H$ is a compact subgroup containing the open
subgroup $K$; hence $H$ is compact open.  Both $C$ and $K$ lie in
$\ker\eta$, and consequently $u\in H$ and $\eta\in H^\perp$.  This proves
(4.1) modulo the two null sets obtained from Tonelli.

Each $H\times H^\perp$ is open in the second-countable space
$G\times\Ghat$.  Their union is therefore an open Lindel\"of subspace, so it
has a countable subcover.  Together with the null-set statement, this proves
the final assertion.
\end{proof}

\section{Pure states give every local Fourier mode}

\begin{lemma}[Fourier basis on one rectangle]\label{lem:local}
For a compact open $H\leq G$,
\[
 \X^0_{P_{H,a,\alpha}}(u,\eta)
 =\one_H(u)\one_{H^\perp}(\eta)\,
   \eta(a)\overline{\alpha(u)}.                              \tag{5.1}
\]
As $(a,\alpha)$ runs through
$G/H\times\Ghat/H^\perp$, these functions form an orthonormal basis of
$L^2(H\times H^\perp)$, extended by zero outside the rectangle.
\end{lemma}

\begin{proof}
Directly from the normalized indicator $e_H$,
\begin{align*}
 \X^0_{P_H}(u,\eta)
 &=\langle e_H,M_\eta T_u e_H\rangle\\
 &=\one_H(u)\one_{H^\perp}(\eta).                            \tag{5.2}
\end{align*}
Indeed, if $u\notin H$ the two cosets are disjoint, while for $u\in H$ the
integral of $\eta$ over $H$ vanishes unless $\eta\in H^\perp$.  Conjugating
$M_\eta T_u$ by $U(a,\alpha)=M_\alpha T_a$ gives
\[
 U(a,\alpha)^*U(u,\eta)U(a,\alpha)
 =\eta(a)\overline{\alpha(u)}U(u,\eta),
\]
which proves (5.1).

The dual of the compact group $H\times H^\perp$ is
\[
 \widehat H\times\widehat{H^\perp}
 \cong (\Ghat/H^\perp)\times(G/H).
\]
Formula (5.1) realizes every character in this dual group.  Fourier theory
on compact abelian groups therefore makes these functions an orthogonal
basis.  Finally,
$\mu_G(H)\mu_{\Ghat}(H^\perp)=1$, so every basis vector has norm one.
\end{proof}

\section{Full resolution of Conjecture (11)}

\begin{theorem}\label{thm:main}
For every second-countable locally compact abelian group $G$,
\[
 \boxed{\displaystyle
 \VKD=\overline{\operatorname{span}_{\mathbb R}(\Epp)}^{\,\|\cdot\|_2}.}
\]
\end{theorem}

\begin{proof}
If $G$ has no compact open subgroup, Theorems 1.2 and 4.1 of \cite{Spriet}
give $\VKD=\{0\}$ and $\Epp=\varnothing$, so the equality is immediate.

Assume that $G$ has a compact open subgroup.  Let $(H_n)_{n\geq1}$ be the
countable family from Lemma~\ref{lem:cover} and set
$R_n=H_n\times H_n^\perp$.  If $f\in L^2(N)$, partition the union into
\[
 E_1=R_1,\qquad E_n=R_n\setminus\bigcup_{j<n}R_j.
\]
Then
$f=\lim_{m\to\infty}\sum_{n\leq m}f\one_{E_n}$ in $L^2$, and each summand
belongs to $L^2(R_n)$.  Lemma~\ref{lem:local} shows that each such summand is
an $L^2$ limit of complex linear combinations of characteristic functions
of elements of $\Epp$.  Hence
\[
 L^2(N)=\overline{\operatorname{span}_{\mathbb C}
   \{\X_P^0:P\in\Epp\}}.                                   \tag{6.1}
\]

Now take $A\in\VKD$.  The source's Lemma 4.1 places $\X_A^0$ in $L^2(N)$.
By (2.1) and (6.1), there are finite complex combinations
$B_m=\sum_jc_{m,j}P_{m,j}$, with $P_{m,j}\in\Epp$, such that
$B_m\to A$ in Hilbert--Schmidt norm.  Since $A=A^*$ and every $P_{m,j}$ is
self-adjoint,
\[
 C_m=\frac{B_m+B_m^*}{2}
     =\sum_j\operatorname{Re}(c_{m,j})P_{m,j}
\]
belongs to $\operatorname{span}_{\mathbb R}(\Epp)$ and
\[
 \|C_m-A\|_2
 \leq\tfrac12\bigl(\|B_m-A\|_2+\|B_m^*-A\|_2\bigr)
 =\|B_m-A\|_2\longrightarrow0.
\]
This proves the nontrivial inclusion.  The reverse inclusion follows because
$\Epp\subseteq\VKD$ and $\VKD$ is a closed real subspace, as established in
the source.
\end{proof}

\section{Verification, novelty, and boundary}

The included verifier treats eight finite abelian groups, enumerates every
subgroup $H$, builds every function in (5.1), and checks both the exact cover
of $N$ and the numerical rank of their span.  In every case the rank equals
$|N|$; this includes the noncyclic groups $(\mathbb Z/2)^2$ and
$(\mathbb Z/2)^3$ and the composite cyclic group $\mathbb Z/6$.  These tests
agree with the finite-group theorem of De Bi\`evre--Langrenez--Radchenko
\cite{DLR} and stress-test the phase factors and quotient parameters.

\begin{center}
\small
\begin{tabular}{c|rrrrrrrr}
$G$ & $\mathbb Z_2$ & $\mathbb Z_3$ & $\mathbb Z_4$ & $\mathbb Z_5$ &
$\mathbb Z_6$ & $\mathbb Z_2^2$ & $\mathbb Z_2\!\times\!\mathbb Z_3$ &
$\mathbb Z_2^3$\\ \hline
$|N|$ &3&5&8&9&15&10&15&36\\
computed rank &3&5&8&9&15&10&15&36
\end{tabular}
\end{center}

The proof's noncomputational checkpoints are: the Haar-positive-subgroup
criterion, compact/open annihilator duality, the Tonelli null-set reduction,
the compactness and openness of $C+K$, the Lindel\"of subcover, and the full
character parametrization in Lemma~\ref{lem:local}.  No positivity claim is
made for a general complex combination; positivity is unnecessary because
Conjecture (11) concerns a real linear space.

Bounded searches of the run indexes, the exact conjecture, the source's later
citations, and KD-real observable terminology through 13 August 2026 found no
later proof for SCLCA groups.  The finite-group analogue is known
\cite{DLR}; the measure-theoretic rectangle cover is the new step claimed
here.  Novelty confidence is moderate-to-high, subject to expert literature
review.

The distinct equality
\[
 \mathcal E_{\mathrm{KD}+}
 =\overline{\operatorname{conv}}(\Epp)
\]
is not proved here.  Attempts to extend the finite filtration decomposition
to an infinite chain of compact open subgroups encounter an uncontrolled
infinite positive-decomposition step.  The present theorem settles exactly
Conjecture (11), not Conjecture (12) or the proposed group characterization.

\noindent\textbf{Human-review recommendation.}
Check the Tonelli reduction in Lemma~\ref{lem:cover}, the annihilator
identifications in Lemma~\ref{lem:local}, and the phase in (5.1).  Then compare
the rectangle-cover argument with terminology for topologically periodic
elements in the LCA-group literature before asserting priority.

\begin{thebibliography}{9}
\bibitem{Spriet}
M. Spriet,
\emph{Characterizing the Kirkwood--Dirac positivity on second countable LCA
groups}, arXiv:2507.23628 (2025; current PDF dated 23 February 2026).

\bibitem{DLR}
S. De Bi\`evre, C. Langrenez, and D. Radchenko,
\emph{The Kirkwood--Dirac representation associated to the Fourier transform
for finite abelian groups: positivity}, Ann. Henri Poincar\'e (2025),
doi:10.1007/s00023-025-01614-7.
\end{thebibliography}

\end{document}
